%%%%%%%%%%  scan idx 57  |  folio 44  %%%%%%%%%%
alors $f(A)=\sum_{i\geq 0}(-1)^i\gamma_{\cC_0}(L^i)$.
On montre que l'élément de $\underline{K}(\cC_0)$ ainsi construit ne dépend pas de la résolution $\underline{L}$ choisie~: si on a une deuxième résolution $\underline{L'}$, on «~coiffe~» $\underline{L}$ et $\underline{L'}$ par une troisième résolution $\underline{L''}$, avec des homomorphismes \emph{surjectifs} de $\underline{L''}$ dans $\underline{L}$ et $\underline{L'}$ (ce qui est possible grâce à (i) et (ii)~; on est alors ramené à prouver $\underline{E}_{\cC_0}(\underline{L})=\underline{E}_{\cC_0}(\underline{L''})$. Or la différence des deux est $\underline{E}_{\cC_0}(\underline{N})$ où $\underline{N}$ est le noyau de l'homomorphisme surjectif de $\underline{L''}$ sur $\underline{L}$, noyau qui est dans $\cC_0$ d'après (i). Or $\underline{N}$ est un complexe acyclique en toutes dimensions, d'où $\underline{E}_{\cC_0}(\underline{N})=0$, d'où la conclusion (N.B. $\underline{E}_{\cC_0}(\underline{L})$, pour un complexe fini à valeurs dans $\cC_0$, désigne la caractéristique d'Euler--Poincaré $\sum_{i\geq0}(-1)^i\gamma_{\cC_0}(L^i)$ de $\underline{L}$).

On prouve de façon analogue que la fonction $f:\cC_1\longrightarrow\underline{K}(\cC_0)$ ainsi construite est \emph{additive}, donc définit un homomorphisme de $\underline{K}(\cC_1)$ dans $\underline{K}(\cC_0)$, et il est immédiat que cet homomorphisme est inverse de l'homomorphisme canonique de $\underline{K}(\cC_0)$ dans $\underline{K}(\cC_1)$. cqfd.

\textbf{Remarque~:} Pour vérifier (ii), il suffit de vérifier que $\cC_1$ est contenu dans une sous-classe $\cC_2$ de $\cC$ satisfaisant aux deux conditions
\begin{enumerate}[label=(ii \alph*)]
\item Tout $A\in\cC_2$ est isomorphe à un quotient d'un $L\in\cC_0$.
\item Si $u:A\to B$, avec $A,B\in\cC_2$, alors $\Ker u\in\cC_2$.
\end{enumerate}

\textbf{Corollaire~:}
\emph{Soit $X$ un espace algébrique quasi-projectif. Soient $\cF(X)$ la catégorie des faisceaux algébriques cohérents sur $X$, $\xi(X)$ la catégorie des faisceaux algébriques cohérents localement libres} (équivalente à la catégorie des fibrés vectoriels sur $X$), \emph{$\cF_0(X)$ la catégorie des faisceaux algébriques cohérents $F$ tels que pour tout $x\in X$, $F_x$ soit un $\cO_x$-module 
